\section{Implementation Methodology}
\label{sec:methodology}

The implementation follows the mathematical model introduced in
Section~\ref{sec:mathematical-model}. The objective is not to build an
independent program with unrelated data structures, but to translate each
mathematical object into a clear software object. The grid cells, admissible
edges, terminal cells, row clues, column clues, and fixed information are first
stored in a normalized instance representation. From this representation, the
solver constructs the grid graph and then builds the mixed-integer model.

\subsection{Software Architecture}

The Python implementation is organized around four main packages. The
\texttt{core} package contains the data model, parser, graph construction
utilities, ASCII display tools, and independent validator. The \texttt{solver}
package contains the PuLP/CBC implementation of the MILP model
\cite{python,pulp}. The \texttt{generation} package creates feasible random
instances by first constructing a path and then deriving row and column clues.
Finally, the \texttt{ui} package contains both a lightweight Pygame viewer and
a separate playable interface.

This separation mirrors the structure of the model. Instance parsing and graph
construction do not depend on the solver. The validator does not depend on the
MILP model. The user interface does not build mathematical constraints. As a
result, the same normalized instance can be solved, validated, rendered,
benchmarked, or used in an interactive game without duplicating the puzzle
rules in several incompatible forms.

\subsection{Instance Representation and File Format}

The central data object is \texttt{TracksInstance}. It stores the grid
dimensions, terminals, row clues, column clues, forced-used cells,
forced-empty cells, forced edges, and fixed local patterns. The code uses
zero-based coordinates, while the mathematical report uses one-based notation.
This is only an indexing convention: a mathematical cell \((i,j)\) corresponds
to the Python tuple \texttt{(i-1, j-1)}.

The text format used for instances is deliberately simple. Required fields are
\texttt{rows}, \texttt{cols}, \texttt{start}, \texttt{end},
\texttt{row\_clues}, and \texttt{col\_clues}. Optional fields encode fixed
information. The field \texttt{fixed\_used} forces a cell to be occupied, but
does not prescribe its shape. The field \texttt{fixed\_empty} forbids a cell
from being used. The field \texttt{fixed\_edges} forces a connection between
two adjacent cells. The field \texttt{fixed\_patterns} is stronger: it fixes
the exact local directions of the track piece in a cell. During normalization,
terminals are automatically treated as used cells, undirected edges are
canonicalized, and inconsistent hints are rejected.

\subsection{Graph Construction and Structural Validation}

The graph layer constructs the objects used in the mathematical formulation:
the vertex set \(V\), horizontal and vertical edge sets \(E^H\) and \(E^V\),
the complete edge set \(E\), neighbor sets \(N(v)\), incident edge sets
\(\delta(v)\), and the directed arc set \(A\). This layer is deterministic and
pure: it only derives graph objects from the normalized instance.

The independent validator is a second line of defense. After a candidate
solution is produced, either by the solver or by the playable interface, the
validator checks the row and column counts, terminal degrees, internal
degrees, fixed hints, connectedness, absence of branches, and absence of
disconnected cycles. This is important for an optimization project because the
solver output should not be trusted solely on the basis of solver status. A
solution is accepted by the application only when it also satisfies the
structural validator.

\subsection{MILP Translation}

The MILP builder translates the constraint families of
Section~\ref{sec:mathematical-model} directly into PuLP expressions. Cell
variables \(y_v\), edge variables \(x_e\), and flow variables \(f_{uv}\) are
created over the grid graph. Row and column constraints are generated by
summing the corresponding cell variables. Edge-cell consistency constraints
link each \(x_e\) to its endpoint variables. Degree constraints impose the
path structure. Fixed information is translated into variable-fixing
constraints. Finally, flow-capacity and flow-balance constraints impose global
connectedness.

The optimization objective is the constant value \(0\), because the puzzle is
a feasibility problem rather than a cost-minimization problem. The MILP solver
therefore searches for any feasible assignment satisfying all constraints. If
PuLP/CBC reports an optimal or feasible assignment, the selected cells and
edges are decoded into a \texttt{TracksSolution} object and passed to the
validator.

\subsection{Constructive Instance Generation}

The generator avoids producing arbitrary clue systems. Instead, it first
constructs a simple path between the terminals. Once this path has been found,
the row and column clues are computed from the cells visited by the path. This
guarantees that every generated instance has at least one feasible solution.
Optional hints are then sampled from the path as forced-used cells or forced
edges.

This procedure does not guarantee uniqueness, and uniqueness is not required
for the Tracks part of the project. Its purpose is to provide controlled
feasible instances for testing, benchmarking, and interactive play. Larger
instances can be generated by increasing the grid size or by requiring longer
paths.

\subsection{Course-Compatible Workflow}

The course reference workflow is expressed through functions such as
\texttt{readInputFile}, \texttt{cplexSolve}, \texttt{displayGrid},
\texttt{displaySolution}, \texttt{generateDataSet}, \texttt{solveDataSet}, and
\texttt{resultsArray}. The project is implemented in Python rather than
Julia, but equivalent wrappers are provided in
\texttt{tracks\_solver.course\_workflow}. The name \texttt{cplexSolve} is kept
as a course-compatible alias even though the backend used here is PuLP/CBC.

The workflow can therefore read a manual instance, display it, solve it with
the MILP backend, display the solution, generate random datasets, solve a
folder of instances, and export result files containing \texttt{solveTime} and
\texttt{isOptimal}. These artifacts make the project easier to compare with
the expected course structure while preserving the Python architecture.

\subsection{Visualization and Playable Interface}

The Pygame layer is separated from the solver. The basic viewer displays an
instance and, when available, an exact solution returned by the MILP backend.
The playable interface adds a menu, random-game generation, map loading,
player editing, explicit validation, and solution reveal. Fixed hints are
rendered with different backgrounds: blocked cells, required cells, forced
connections, and exact fixed patterns are visually distinguished.

The UI is not required for the mathematical correctness of the solver, but it
is useful for interpreting instances and for defense demonstrations. In
particular, it makes visible the difference between local track placement and
global graph validity: a visually plausible partial route may still fail the
degree, clue, or connectivity checks.
